\section{Operational evaluation and prospective collection}\label{sec:section-09}

The active r2 runtime passed 27 source-bound tests. All 8 actual process-kill checkpoints recovered with equivalent state, and 5 altered-release evidence cases were refused. The operating system invoked the frozen command at 2026-09-28T08:50:56.358915+00:00; its parent PID was 1. That commissioning run validated 22 instruments, prepared signals and targets, persisted both accounts and reconciled them. Its receipt was then restored at the registered original path while preserving the prior directory. Commissioning is isolated from the pilot: the verified observed decision count is 0, mature count 0, and terminal mark count 0. The schedule's next event is 2026-09-30T04:30:00+08:00 Hong Kong. No remote trading workflow was executed; the inspected GitHub workflow was disabled by inactivity. Physical sleep/power-loss and sustained session-close availability remain unobserved.

The actual cold provider-to-ledger commissioning took 33.27 seconds, including 30.11 seconds of public-data reads. It made 22 requests with 0 observed provider failures. Across 20 warm batches, median batch time was 1.798 seconds and empirical P95 was 4.159 seconds. A separate 5-fit SPY timing sample had median 0.424 seconds. Cold-run and small training samples do not support tail-latency estimates; warm P95 is descriptive. The collector ran on the recorded arm64 macOS host under concurrent research load.

\noindent Table~\ref{tab:17} reports warm pipeline stage timings, seconds; P99 intentionally unreported.

\begin{table}[!htbp]\centering\footnotesize\setlength{\tabcolsep}{3pt}\renewcommand{\arraystretch}{1.14}
\caption{Warm pipeline stage timings, seconds; P99 intentionally unreported}\label{tab:17}
\begin{tabular}{@{}>{\raggedright\arraybackslash}p{\dimexpr 0.22000\textwidth-4.800pt\relax}>{\raggedleft\arraybackslash}p{\dimexpr 0.19500\textwidth-4.800pt\relax}>{\raggedleft\arraybackslash}p{\dimexpr 0.19500\textwidth-4.800pt\relax}>{\raggedleft\arraybackslash}p{\dimexpr 0.19500\textwidth-4.800pt\relax}>{\raggedleft\arraybackslash}p{\dimexpr 0.19500\textwidth-4.800pt\relax}@{}}
\toprule
\textbf{Stage} & \textbf{Samples} & \textbf{P50 (s)} & \textbf{P95 (s)} & \textbf{P99 (s)} \\
\midrule
features & 20 & 0.372 & 1.083 & - \\
inference & 20 & 0.191 & 0.354 & - \\
whole\_\allowbreak{}batch & 20 & 1.798 & 4.159 & - \\
\bottomrule\end{tabular}\par\smallskip
\raggedright\textit{Source:} \path{data/research-v2/operations/latency-r2.json}.
\end{table}


\noindent Table~\ref{tab:18} reports actual verification jobs; overlapping suites are separate records.

\begin{table}[!htbp]\centering\footnotesize\setlength{\tabcolsep}{3pt}\renewcommand{\arraystretch}{1.14}
\caption{Actual verification jobs; overlapping suites are separate records}\label{tab:18}
\begin{tabular}{@{}>{\raggedright\arraybackslash}p{\dimexpr 0.56180\textwidth-4.800pt\relax}>{\raggedleft\arraybackslash}p{\dimexpr 0.10955\textwidth-4.800pt\relax}>{\raggedleft\arraybackslash}p{\dimexpr 0.10955\textwidth-4.800pt\relax}>{\raggedleft\arraybackslash}p{\dimexpr 0.10955\textwidth-4.800pt\relax}>{\raggedleft\arraybackslash}p{\dimexpr 0.10955\textwidth-4.800pt\relax}@{}}
\toprule
\textbf{Saved report} & \textbf{Cases} & \textbf{Failed} & \textbf{Errors} & \textbf{Skipped} \\
\midrule
data/\allowbreak{}research-v2/\allowbreak{}tests/\allowbreak{}baseline-offline.xml & 310 & 0 & 0 & 0 \\
data/\allowbreak{}research-v2/\allowbreak{}tests/\allowbreak{}final-offline/\allowbreak{}results.xml & 370 & 0 & 0 & 0 \\
data/\allowbreak{}research-v2/\allowbreak{}tests/\allowbreak{}final-qualified/\allowbreak{}results.xml & 370 & 0 & 0 & 0 \\
data/\allowbreak{}research-v2/\allowbreak{}tests/\allowbreak{}final-stable/\allowbreak{}results.xml & 370 & 0 & 0 & 0 \\
data/\allowbreak{}research-v2/\allowbreak{}tests/\allowbreak{}provider-integration.xml & 17 & 0 & 0 & 0 \\
data/\allowbreak{}research-v2/\allowbreak{}tests/\allowbreak{}ui-integration.xml & 1 & 0 & 0 & 0 \\
data/\allowbreak{}research-v2/\allowbreak{}tests/\allowbreak{}ui-r2/\allowbreak{}results.xml & 1 & 0 & 0 & 0 \\
data/\allowbreak{}research-v2/\allowbreak{}tests/\allowbreak{}ui-r2-qualified/\allowbreak{}results.xml & 1 & 0 & 0 & 0 \\
\bottomrule\end{tabular}\par\smallskip
\raggedright\textit{Source:} \path{data/research-v2/tests/*.xml; XML testcase counts}.
\end{table}


\noindent Figure~\ref{fig:13} shows measured fitting latency by registered variant under differing run order and concurrent host load; these are operational timings, not controlled speedup estimates.

\begin{figure}[htbp]\centering
% Native editable figure; frozen statistics 0d291b4aebf8c524dc6705d872c070feb990e86a1cebbd59a758f400b240c952
\begin{tikzpicture}
\begin{axis}[width=.9\linewidth,height=6cm,boxplot/draw direction=y,xtick={0,1,2,3,4,5,6,7,8,9},xticklabels={{base},{cadence\_21},{fraction\_0.5},{fraction\_0.75},{history\_504},{transfer},{without\_context\_volume},{without\_momentum},{without\_risk},{without\_trend}},x tick label style={rotate=35,anchor=east},ylabel={Fit seconds},ymin=0]
\addplot+[boxplot prepared={draw position=0,lower whisker=0.187755082996,lower quartile=0.230702500245,median=0.667628895499,upper quartile=0.966208124744,upper whisker=1.733549875},blue] coordinates {};
\addplot+[boxplot prepared={draw position=1,lower whisker=0.187187374991,lower quartile=0.192843219254,median=0.197883166991,upper quartile=0.201251249491,upper whisker=0.212719375006},blue] coordinates {};
\addplot+[boxplot prepared={draw position=2,lower whisker=0.145877624993,lower quartile=0.152505896251,median=0.155188041506,upper quartile=0.158791937243,upper whisker=0.166472625002},blue] coordinates {};
\addplot+[boxplot prepared={draw position=3,lower whisker=0.174998332994,lower quartile=0.181211584,median=0.184082833497,upper quartile=0.185834291995,upper whisker=0.191644124992},blue] coordinates {};
\addplot+[boxplot prepared={draw position=4,lower whisker=0.133230000007,lower quartile=0.137934395501,median=0.1396283125,upper quartile=0.147123426996,upper whisker=0.157956458992},blue] coordinates {};
\addplot+[boxplot prepared={draw position=5,lower whisker=0.779778124997,lower quartile=0.983922364248,median=1.07975004199,upper quartile=1.443187437,upper whisker=1.45833337501},blue] coordinates {};
\addplot+[boxplot prepared={draw position=6,lower whisker=0.174084333994,lower quartile=0.18489948925,median=0.241868416495,upper quartile=0.318596260506,upper whisker=0.500081375008},blue] coordinates {};
\addplot+[boxplot prepared={draw position=7,lower whisker=0.186025834002,lower quartile=0.22113709425,median=0.275304978997,upper quartile=0.392184739751,upper whisker=0.581473124999},blue] coordinates {};
\addplot+[boxplot prepared={draw position=8,lower whisker=0.174779500012,lower quartile=0.207700540999,median=0.244683541998,upper quartile=0.338147614493,upper whisker=0.484661540992},blue] coordinates {};
\addplot+[boxplot prepared={draw position=9,lower whisker=0.17080854201,lower quartile=0.185769979245,median=0.199566438001,upper quartile=0.216780635245,upper whisker=0.258316874999},blue] coordinates {};

\end{axis}
\end{tikzpicture}

\caption{Measured fitting latency by registered variant under differing run order and concurrent host load; these are operational timings, not controlled speedup estimates.}\label{fig:13}
\end{figure}

The local scheduler uses exchange-session events rather than a fixed Hong Kong clock. Standard close+30 occurs at different local times across US daylight-saving changes, and early-close sessions advance the event. Only an eligible transient failure permits the registered retries; deterministic integrity or risk failures halt. The watchdog checks for an absent receipt as well as an explicitly failed one. Pre-open checks do not create replacement strategy decisions. Weekly archival and milestone reviews assess coverage, latency, risk and protocol deviations without choosing winners from pilot returns.

The collector cannot execute while its host is asleep or unavailable. On recovery it must distinguish an overdue event from a legitimate current event, record missed coverage and refuse backfilled prospective decisions. A real launchd trigger receipt demonstrates that the operating system invoked the current runner and it reached its persistent state. It does not demonstrate a future session-close invocation, sustained host availability or an always-on service. Those are observations that must accumulate.

\noindent Table~\ref{tab:19} reports registered XNYS events and Hong Kong display times.

\begin{table}[!htbp]\centering\footnotesize\setlength{\tabcolsep}{3pt}\renewcommand{\arraystretch}{1.14}
\caption{Registered XNYS events and Hong Kong display times}\label{tab:19}
\begin{tabular}{@{}>{\raggedright\arraybackslash}p{\dimexpr 0.49153\textwidth-4.000pt\relax}>{\raggedleft\arraybackslash}p{\dimexpr 0.16949\textwidth-4.000pt\relax}>{\raggedright\arraybackslash}p{\dimexpr 0.33898\textwidth-4.000pt\relax}@{}}
\toprule
\textbf{Event} & \textbf{US session} & \textbf{Hong Kong trigger} \\
\midrule
Observed-session milestone 1 & 2026-09-29 & 2026-09-30 04:30 HKT \\
Observed-session milestone 5 & 2026-10-05 & 2026-10-06 04:30 HKT \\
Observed-session milestone 21 & 2026-10-27 & 2026-10-28 04:30 HKT \\
Observed-session milestone 42 & 2026-11-25 & 2026-11-26 05:30 HKT \\
Observed-session milestone 63 & 2026-12-28 & 2026-12-29 05:30 HKT \\
Model refit at anchor+21 & 2026-10-28 & 2026-10-29 04:30 HKT \\
Model refit at anchor+42 & 2026-11-27 & 2026-11-28 02:30 HKT \\
Final target-label endpoint & 2027-01-06 & 2027-01-07 05:30 HKT \\
\bottomrule\end{tabular}\par\smallskip
\raggedright\textit{Source:} \path{data/research-v2/tables/schedule_milestones.csv}.
\end{table}


The latest validated prospective analysis contains 0 decision receipts and 0 mature decision dates. It is recorded as interim. Commissioning, controlled replay and historical predictions are excluded from these counts. No prospective forecast or economic result is available; unobserved quantities are not zero returns or successful sessions.

Mature outcomes attach to earlier immutable decisions only after the execution-aligned endpoint is observed. After the final decision, six terminal sessions process outstanding fills and marks without generating new allocation targets. A halted account can still retain valid matured labels, but missing terminal accounting prevents a complete economic freeze. The read-only final-report command reconstructs the costed equal-weight comparator from those same observed snapshots and releases inferential results only at the registered final freeze; interim views are descriptive. The last decision date and the last label-maturity date are different. A missing expected session remains a recorded outage and does not shift the registered end until results become favorable. The operational observation period is not a promise of adequate power for a small economic effect. A strategy can complete the operational pilot and still fail the statistical or economic gate.
